% संपूर्ण मराठी ग्रंथातील प्रकरण 57: किमान बदलाची चिन्हार्थमीमांसा
% हा संपादनयोग्य स्रोतखंड आहे. संकलनासाठी संपूर्ण स्रोतसंग्रहातील
% अक्षररूपे, पूर्वभाग, चिन्हव्याख्या आणि आकृती-स्रोत आवश्यक आहेत.
% मूळ: Open Logic Project; मजकूर आणि रूपांतर CC BY 4.0.
% यंत्रानुवाद, दुरुस्ती आणि तपासणी: OpenAI Codex —
% GPT-5.6 Sol आणि GPT-6 Sol, दोन्ही Ultra effort.
% दुरुस्ती-पुनरावलोकन: GPT-6.1 Sol, Ultra effort.
\chapter{किमान बदलाची चिन्हार्थमीमांसा}\label{cnt:min:chap}










\section{प्रस्तावना}\label{cnt:min:int:sec}

‘‘ही आगकाडी ओढली असती, तर ती पेटली असती’’
अशा प्रतिवास्तविक सशर्त विधानांचे
स्टालनाकर आणि लुईस यांनी विश्लेषण सुचवले.
अशा वाक्यांच्या सत्यतेच्या अटी
योग्य रीतीने कशा समजाव्यात,
याविषयी त्यांची मांडणी होती.
दोन्ही मांडण्यांमागील कल्पना अशी आहे:
प्रतिवास्तविक सशर्त विधान सत्य आहे का
हे ठरवताना त्याचे पूर्वांग सत्य होण्यासाठी
वास्तविक जगापासून किमान फरक असलेली
संभाव्य जगे विचारात घ्यावी लागतात.
या संभाव्य जगांत उत्तरांग सत्य असेल,
तर प्रतिवास्तविक विधान सत्य असते.
उदाहरणार्थ, माझ्या हातात आगकाडी
आणि आगपेटी आहे असे समजा.
वास्तविक जगात मी फक्त त्यांच्याकडे
पाहतो आणि आगकाडी ओढली तर काय
होईल, याचा विचार करतो.
मी आगकाडी ओढतो अशा पर्यायी जगातला
किमान बदल म्हणजे मी कृती करण्याचे
ठरवतो आणि आगकाडी ओढतो.
हा बदल किमान आहे, कारण आगकाडी ओढणे आणि त्याचे
आवश्यक परिणाम यांखेरीज दुसरे अनावश्यक बदल होत नाहीत: मी त्याच वेळी
हवेत उडी मारत नाही; आगकाडी
ओढल्याने माझ्या केसांनाही आग
लागत नाही; माझ्या बोटांतील सारी
ताकद अचानक नाहीशी होत नाही;
त्याच वेळी कुणी सुपर-सोकरने
माझ्यावर पाणी मारत नाही;
इत्यादी. त्या पर्यायी शक्यतेत
आगकाडी पेटते. म्हणून मी
आगकाडी ओढली असती,
तर ती पेटली असती,
हे सत्य आहे.

या सहज समजणाऱ्या मांडणीला
प्रतिवास्तविक तर्कशास्त्रांची औपचारिक
चिन्हार्थमीमांसा जोडता येते.
प्रतिवास्तविक विधानासाठी लुईस यांनी
‘‘$\boxright$’’ हे चिन्ह, तर
स्टालनाकर यांनी ‘‘$>$’’ हे चिन्ह
वापरले. आपण~$\cif$ वापरू आणि
ते विधानीय तर्कशास्त्रात द्विपदी
संयोजक म्हणून जोडू.
म्हणून $\varphi \lif \psi$ या रूपाच्या
सूत्रे बरोबरच
$\varphi \cif \psi$ या रूपाची
सूत्रे ही असतील.
मोडल तर्कशास्त्राच्या संबंधात्मक
चिन्हार्थमीमांसेप्रमाणे, औपचारिक
चिन्हार्थमीमांसा अशा प्रतिमानांवर
आधारलेली असते की ज्यांत
सूत्रे यांचे मूल्यांकन जगांच्या
संदर्भात केले जाते. त्यात
$\mSat{M}{\varphi \cif \psi}[w]$
याची पूर्ति-अट काही (इतर) जगे~$w'$
यांच्या संदर्भातील
$\mSat{M}{\varphi}[w']$ आणि
$\mSat{M}{\psi}[w']$
यांच्या साहाय्याने दिली जाते.
कोणती $w'$ जगे?
सहज कल्पनेनुसार,
$\mSat{M}{\varphi}[w']$ सत्य असलेली
आणि~$w$ च्या सर्वाधिक जवळ
असलेली जगे. त्यासाठी
‘‘जवळीक’’ दर्शवणारा
संबंधही प्रतिमानात
समाविष्ट करावा लागतो.

प्रतिवास्तविक परिस्थिती चित्ररूपात
दाखवण्याची एक बोधक पद्धत लुईस
यांनी आणली. प्रत्येक संभाव्य जग
हे इतर जगे समाविष्ट करणाऱ्या
एकमेकांच्या आत असलेल्या गोलकांच्या
प्रणालीच्या केंद्रस्थानी असते;
आपण हे गोलक समकेंद्री
वर्तुळांनी दाखवतो. या आकृतीत शेजारच्या दोन गोलकांच्या मधला थर
समान जवळीक दाखवतो. अधिक सूक्ष्म गोलकांची सामान्य
प्रणाली असताना मनमानी दोन गोलकांमधील सर्व जगे
सारख्याच जवळ असण्याची अट नसते.
आतल्या गोलकात असलेली जगे
अधिक जवळ, तर बाहेरच्या
गोलकात असलेली जगे अधिक
दूर असतात.
\begin{center}
\begin{tikzpicture}[scale=.7]\small
  \spheresystem{5}
  \draw (0,0) node {$w$};
  \propositionintersect{0}{4}{40}{3.5}
  {\tikzset{proposition/.append={smooth,tension=1.4}}
  \proposition[shift={(1.6,-1)}]{90}{1}{30}{4}}
  \path (1.5,3) node[below] {$\formula{B}$};
  \path (3,0) node {$\formula{A}$};
\end{tikzpicture}
\end{center}
या आकृतीतील सर्वाधिक जवळची $\varphi$-जगे म्हणजे
$\varphi$ सत्य असलेले किमान एक जग धारण करणाऱ्या
सर्वांत लहान गोलकातील $\varphi$-जगे~$w'$ (राखाडी भाग).
तो प्रणालीतील सर्वांत लहान गोलकच असावा असे नाही.
असा $\varphi$-जग धारण करणारा सर्वांत लहान गोलक
सर्व प्रतिमानांत अस्तित्वात असतोच असे नाही;
पुढील विभागातील गोलक-अट अशा प्रकरणांनाही लागू होते.
सहज कल्पनेनुसार,
सर्वाधिक जवळच्या सर्व
$\varphi$-जगांत $\psi$ सत्य
असेल, तर~$w$ येथे
$\varphi \cif \psi$ ची
पूर्ति होते.













\section{गोलक प्रतिमाने}\label{con:min:sph:sec}

प्रतिवास्तविक विधानांना औपचारिक
चिन्हार्थमीमांसा देण्याचा एक मार्ग
म्हणजे लुईस यांच्या अनौपचारिक
मांडणीचे गणिती रचनेत रूपांतर करणे.
मग जग~$w$ भोवतीचे गोलक
हे जगांचे संच असतात.
गोलक एकमेकांच्या आत
असल्यामुळे, जग~$w$ भोवतीच्या
या जगांच्या संचांना उपसंच-संबंधाने
रेषीय क्रमात असावे लागते.

\begin{defn}
  \emph{गोलक प्रतिमान} म्हणजे
  $\mModel{M} = \tuple{W, O, V}$
  हे त्रिक: येथे $W$ हा जगांचा
  रिक्त नसलेला संच आहे,
  $V\colon \PVar \to \Pow{W}$
  हे सत्य-मूल्यांकन आहे, आणि
  $O\colon W \to \Pow{\Pow{W}}$
  प्रत्येक जग~$w$ ला
  \emph{गोलकांची प्रणाली}~$O_w$
  नेमून देते. प्रत्येक~$w$ साठी
  $O_w$ हा जगांच्या संचांचा संच
  असून त्याने पुढील अटी पूर्ण
  केल्या पाहिजेत:
  \begin{enumerate}
  \item $O_w$ चे \emph{केंद्र}~$w$
    असते: $\{w\} \in O_w$.
  \item $O_w$ मधील गोलक
    \emph{एकमेकांच्या आत} असतात:
    $S_1$, $S_2 \in O_w$ असतील,
    तर $S_1 \subseteq S_2$ किंवा
    $S_2 \subseteq S_1$; म्हणजे
    $O_w$ ला~$\subseteq$ ने
    रेषीय क्रम असतो.
  \item प्रत्येक रिक्त नसलेल्या उपकुल $\mathcal S\subseteq O_w$
    साठी $\bigcup\mathcal S\in O_w$.
  \item प्रत्येक रिक्त नसलेल्या उपकुल $\mathcal S\subseteq O_w$
    साठी $\bigcap\mathcal S\in O_w$.
    या अटी अनंत उपकुलांनाही लागू होतात.
  \end{enumerate}
\end{defn}

$O_w$ मागील कल्पना अशी:
$w$ भोवतीची जगे~$w$ पासून
किती दूर आहेत यानुसार
थरांत मांडली जातात.
रिकामे गोलक वगळल्यास सर्वांत आतला गोलक म्हणजे
एकटे~$w$, अर्थात~$\{w\}$: इतर कोणत्याही जगापेक्षा
$w$ हे स्वतःच्या अधिक जवळ असते. व्याख्या रिकामा
गोलकही परवानगी देते; त्याने पुढील सत्यता-अट बदलत नाही.
$S \subsetneq S'$ असेल,
तर $S' \setminus S$
मधील जगे~$S$ मधील
जगांपेक्षा~$w$ पासून
अधिक दूर असतात:
$S'\setminus S$ हा $S'$ च्या आत पण $S$ च्या बाहेर
असलेल्या जगांचा ‘‘थर’’ आहे.
विशेषतः, गोलकांत
त्यांच्या बाह्य पृष्ठभागाच्या
आतील सर्व जगे असतात
असे समजावे; गोलक
म्हणजे फक्त स्वतंत्र
थर नव्हेत.

\begin{figure}
\begin{center}
\begin{tikzpicture}[layerwidth=1.5,scale=.6]\tiny
  \spheresystem{3}
  \propositionintersect{0}{3}{60}{6}
  \path (0,0) node[world] {$w$};
  \spherepos{-30}{2}{node[world] {$w_2$}}
  \spherepos{220}{2}{node[world] {$w_3$}}
  \spherepos{90}{2}{node[world] {$w_1$}}
  \spherepos[shift={(.1,0)}]{10}{3}{node[world] {$w_5$}}
  \spherepos[shift={(.1,0)}]{-10}{3}{node[world] {$w_6$}}
  \spherepos{225}{3}{node[world] {$w_4$}}
  \spherepos{20}{4}{node[world] {$w_7$}}
  \path (5,0) node[left] {$p$};
\end{tikzpicture}
\caption{गोलक प्रतिमानाची आकृती}
\label{con:min:sph:fig:sphere-model}
\end{center}
\end{figure}

\ref{con:min:sph:fig:sphere-model} मधील
आकृती पुढील गोलक प्रतिमान
दाखवते: $W = \{w, w_1, \dots, w_7\}$,
$V(p) = \{w_5, w_6,
w_7\}$. सर्वांत आतला गोलक
$S_1 = \{w\}$ आहे.
$w$ च्या सर्वाधिक जवळची
जगे $w_1, w_2, w_3$
आहेत; म्हणून पुढचा मोठा
गोलक $S_2 = \{w, w_1,
w_2, w_3\}$ आहे.
त्याहून बाहेरची जगे
$w_4$, $w_5$, $w_6$
आहेत; म्हणून सर्वांत
बाहेरचा गोलक
$S_3 = \{w, w_1, \dots, w_6\}$
आहे. $w$ भोवतीची
गोलक-प्रणाली
$O_w = \{S_1, S_2, S_3\}$
आहे. जग~$w_7$ हे~$w$
भोवतीच्या कोणत्याही गोलकात
नाही. ज्या जगांत~$p$
सत्य आहे त्यांपैकी
सर्वाधिक जवळची जगे
$w_5$ आणि~$w_6$
आहेत. त्यामुळे~$p$
सत्य असलेले जग
समाविष्ट करणारा सर्वांत
लहान गोलक~$S_3$ आहे.

गोलक प्रतिमान~$\mModel M$
मधील जग~$w$ येथे
सूत्र~$\varphi$ ची पूर्ति,
$\mSat{M}{\varphi}[w]$,
याची व्याख्या करताना
मोडल सूत्रे
यांच्या व्याख्येत
$\psi \cif \chi$
साठी पुढील अट जोडतो:

\begin{defn}
  $\mSat{M}{\psi \cif \chi}[w]$
  तेव्हा आणि केवळ तेव्हाच,
  जेव्हा पुढीलपैकी एक घडते:
  \begin{enumerate}
  \item\label{con:min:sph:sphere-vac}
    सर्व $u \in \bigcup O_w$
    साठी $\mSat/{M}{\psi}[u]$; किंवा
  \item\label{con:min:sph:sphere-nonvac}
    एखाद्या $S \in O_w$ साठी,
    \begin{enumerate}
      \item एखाद्या $u \in
        S$ साठी
        $\mSat{M}{\psi}[u]$; आणि
      \item सर्व $v \in S$
        साठी, एकतर
        $\mSat/{M}{\psi}[v]$
        किंवा
        $\mSat{M}{\chi}[v]$.
    \end{enumerate}
  \end{enumerate}
\end{defn}

या व्याख्येनुसार
$\mSat{M}{\psi \cif \chi}[w]$
तेव्हा आणि केवळ तेव्हाच, जेव्हा
एकतर पूर्वांग~$\psi$
हे~$w$ भोवतीच्या सर्व
गोलकांत सर्वत्र असत्य असते,
किंवा एखाद्या गोलक~$S$
मध्ये~$\psi$ सत्य असते आणि
त्या ‘‘$\psi$ समाविष्ट
करणाऱ्या’’ गोलकातील
सर्व जगांत
$\psi \lif \chi$
हे वास्तविक अभिव्यंजन
सत्य असते.
सहज स्पष्टीकरणावरून
अपेक्षा वाटू शकते,
तरी व्याख्येत~$S$
हा~$\psi$ समाविष्ट
करणारा \emph{सर्वांत
आतला} गोलक असावा
अशी अट घातलेली नाही.
पण \ref{con:min:sph:sphere-nonvac}
मधील अट एखाद्या
गोलक~$S$ साठी
पूर्ण होत असेल,
तर~$S$ मध्ये
समाविष्ट असलेल्या आणि
पूर्वांग सत्य असलेले जग
धारण करणाऱ्या प्रत्येक
लहान गोलकासाठीही
ती पूर्ण होते.
म्हणून असा सर्वांत
आतला गोलक अस्तित्वात
असेल, तर त्यासाठीही
अट पूर्ण होते.

गोलक प्रतिमानाच्या व्याख्येनुसार $\psi$ सत्य असलेले
जग धारण करणारा सर्वांत आतला गोलक असतोच असे नाही.
केवळ अनंत उतरती क्रमिका दिल्याने तो नसल्याचे सिद्ध होत नाही;
प्रणालीत तिच्यापेक्षा लहान पूर्वांग-गोलकही असू शकतो.
एक पूर्ण उदाहरण असे: $W=\{w\}\cup\{u_n:n\geq1\}$,
$S_n=\{w\}\cup\{u_j:j\geq n\}$ आणि
$O_w=\{\{w\}\}\cup\{S_n:n\geq1\}$ घ्या.
इतर जगांसाठी $O_{u_n}=\{\{u_n\}\}$ घ्या आणि
पूर्वांग म्हणून $\psi=p$, $V(p)=\{u_n:n\geq1\}$ ठेवा.
ही प्रणाली केंद्रित, एकमेकांच्या आत असलेली आणि प्रत्येक
रिक्त नसलेल्या उपकुलाच्या संघटनाखाली व छेदाखाली बंद आहे.
गोलक $S_n$ हे सर्व $p$-जग धारण करतात,
$S_{n+1}\subsetneq S_n$ आणि $\bigcap_{n\geq1}S_n=\{w\}$,
पण $p$ हे $w$ येथे असत्य आहे. म्हणून सर्वांत लहान
$p$-जग धारण करणारा गोलक नाही.
या प्रतिमानात $\mSat{M}{p\cif \chi}[w]$ तेव्हा आणि केवळ तेव्हाच,
जेव्हा एखाद्या $i\geq1$ साठी सर्व $j\geq i$ येथे
$\mSat{M}{\chi}[u_j]$. सममूल्यपणे, अशा $i$ पासून
$S_i,S_{i+1},\dots$ या सर्व गोलकांत $p\lif \chi$ सत्य असते.
म्हणून सर्वांत लहान पूर्वांग-गोलक नसतानाही मूळ सत्यता-अट
निर्णायक आहे.













\section{प्रतिवास्तविक विधानांचे सत्य आणि असत्य}\label{con:min:tf:sec}

सर्वाधिक जवळची $\varphi$-जगे अस्तित्वात असतील आणि
ती सर्व $\psi$-जगे असतील, तर प्रतिवास्तविक
विधान $\varphi \cif \psi$
(रिक्त नसलेल्या अर्थाने) सत्य असते.
\ref{con:min:tf:fig:true} मध्ये हे दाखवले आहे.
\begin{figure}
\begin{center}
\begin{tikzpicture}[scale=.7]\small
  \spheresystem{5}
  \draw (0,0) node {$w$};
  \propositionintersect{0}{4}{40}{3.5}
  {\tikzset{proposition/.append={smooth,tension=1.4}}
  \proposition[shift={(1.6,-1)}]{90}{1}{30}{4}}
  \path (1.5,3) node[below] {$\formula{B}$};
  \path (3,0) node {$\formula{A}$};
\end{tikzpicture}
\caption{रिक्त नसलेल्या अर्थाने सत्य प्रतिवास्तविक विधान}
\label{con:min:tf:fig:true}
\end{center}
\end{figure}
जग~$w$ भोवतीच्या गोलकांच्या
प्रणालीत $\varphi$ समाविष्ट
करणारा एकही गोलक नसेल,
तरी~$w$ येथे प्रतिवास्तविक
विधान सत्य असते. अशा वेळी
ते \emph{रिक्तपणे} सत्य
असते (\ref{con:min:tf:fig:vacuous} पाहा).
\begin{figure}
\begin{center}
\begin{tikzpicture}[scale=.7]\small
  \spheresystem{5}
  \draw (0,0) node {$w$};
  \proposition{0}{6.5}{40}{3.5}
  {\tikzset{proposition/.append={smooth,tension=1.4}}
  \proposition[shift={(1.2,-1)}]{90}{1}{30}{4}}
  \path (1.5,3) node[below] {$\formula{B}$};
  \spherepos{0}{7}{node {$\formula{A}$}}
\end{tikzpicture}
\caption{रिक्तपणे सत्य प्रतिवास्तविक विधान}
\label{con:min:tf:fig:vacuous}
\end{center}
\end{figure}

सर्वाधिक जवळची $\varphi$-जगे अस्तित्वात असलेल्या या
आकृत्यांत प्रतिवास्तविक विधानाच्या असत्यतेचे दोन
प्रकार पाहू. ही मांडणी सर्वाधिक जवळचे जग नसलेल्या
प्रतिमानांची पूर्ण विभागणी असल्याचा दावा नाही.
पहिल्या प्रकारात सर्वाधिक
जवळची $\varphi$-जगे
सर्वच $\psi$-जगे नसतात,
पण त्यांपैकी काही
तशी असतात.
या वेळी $\varphi \cif
\lnot \psi$ हे विधानही
असत्य असते
(\ref{con:min:tf:fig:false} पाहा).
\begin{figure}
\begin{center}
\begin{tikzpicture}[scale=.7]\small
  \spheresystem{5}
  \draw (0,0) node {$w$};
  \propositionintersect{0}{4}{40}{3.5}
  {\tikzset{proposition/.append={smooth,tension=1.4}}
  \proposition[shift={(1.6,0)}]{90}{1}{30}{3}}
  \path (1.5,3) node[below] {$\formula{B}$};
  \path (3,0) node {$\formula{A}$};
\end{tikzpicture}
\caption{असत्य प्रतिवास्तविक विधान; विरुद्ध उत्तरांगाचे विधानही असत्य}
\label{con:min:tf:fig:false}
\end{center}
\end{figure}
सर्वाधिक जवळच्या
$\varphi$-जगांचा
$\psi$-जगांशी
अजिबात छेद नसेल,
तर $\varphi \cif \psi$
असत्य असते. पण
या वेळी सर्वाधिक
जवळची सर्व
$\varphi$-जगे
$\lnot \psi$-जगे
असतात; म्हणून
$\varphi \cif \lnot \psi$
सत्य असते
(\ref{con:min:tf:fig:false-opposite} पाहा).
\begin{figure}
\begin{center}
\begin{tikzpicture}[scale=.7]\small
  \spheresystem{5}
  \draw (0,0) node {$w$};
  \propositionintersect{0}{4}{40}{3.5}
  {\tikzset{proposition/.append={smooth,tension=1.4}}
  \proposition[shift={(-1,0)}]{90}{1}{30}{3}}
  \path (-1,3) node[below] {$\formula{B}$};
  \path (3,0) node {$\formula{A}$};
\end{tikzpicture}
\caption{असत्य प्रतिवास्तविक विधान; विरुद्ध उत्तरांगाचे विधान सत्य}
\label{con:min:tf:fig:false-opposite}
\end{center}
\end{figure}

S5 मधील कठोर अभिव्यंजनाच्या सत्यतेच्या अवश्यतेशी
तुलना करता, प्रतिवास्तविक
विधाने परिस्थितीवश
असू शकतात.
\ref{con:min:tf:fig:contingent} मधील
गोलक प्रतिमान पाहा.
$u$ च्या सर्वाधिक
जवळची $\varphi$-जगे
ही सर्व $\psi$-जगे
आहेत; म्हणून
$\mSat{M}{\varphi \cif
  \psi}[u]$. पण~$v$
च्या सर्वाधिक जवळच्या
काही $\varphi$-जगांत
$\psi$ सत्य नाही;
म्हणून
$\mSat/{M}{\varphi \cif \psi}[v]$.
\begin{figure}
\begin{center}
\begin{tikzpicture}[scale=.6]\small
  \spheresystem{7}
  \spheresystem[shift={(0:3.1)},dashed]{4}
  \propositionintersect{45}{4}{40}{4.5}
  \begin{scope}
    \clip \propositionplot{45}{4}{40}{4.5} ;
    \spherelayer[shift={(0:3.1)}]{3}
  \end{scope}
  \proposition[shift={(1.4,-.5)}]{90}{1}{35}{5}
  \path (0,0) node {$u$};
  \path (0:3.1) node {$v$};
  \spherepos{35}{9}{node {$\formula{A}$}}
  \spherepos{80}{9}{node {$\formula{B}$}}
\end{tikzpicture}
\end{center}
\caption{परिस्थितीवश प्रतिवास्तविक विधान}
\label{con:min:tf:fig:contingent}
\end{figure}











\section{पूर्वांग बळकट करणे}\label{cnt:min:agg:sec}

‘‘पूर्वांग बळकट करणे’’ म्हणजे
$\varphi \lif \chi \Entails
(\varphi \land \psi) \lif \chi$
हे अनुमान. वास्तविक
अभिव्यंजनासाठी ते वैध,
पण प्रतिवास्तविक
विधानांसाठी अवैध आहे.
मी ही आगकाडी ओढली असती,
तर ती पेटली असती,
हे सत्य आहे असे समजा.
(म्हणजे आगकाडीत किंवा
आगपेटीच्या घासण्याच्या
पृष्ठभागात काही दोष नाही,
मी आगकाडी मोडणार नाही,
इत्यादी.) पण ही
आगकाडी मी बाह्य
अवकाशात ओढली असती,
तर ती पेटली असती,
हे सत्य नाही. म्हणून
पुढील अनुमान अवैध आहे:
\begin{quote}
  आगकाडी ओढली असती,
  तर ती पेटली असती.

  म्हणून आगकाडी बाह्य
  अवकाशात ओढली असती,
  तर ती पेटली असती.
\end{quote}

लुईस--स्टालनाकर यांच्या
सशर्त विधानांच्या
विश्लेषणातून हे
समजते: मी आगकाडी
बाह्य अवकाशात
ओढतो असे सर्वाधिक
जवळचे जग, मी
आगकाडी ओढतो
अशा सर्वाधिक
जवळच्या जगापेक्षा
वास्तविक जगापासून
खूप दूर आहे.
म्हणून दुसऱ्या
जगात आगकाडी
पेटते हे सत्य
असले, तरी
पहिल्या जगात
ती पेटत नाही.
असेच असायला हवे.

\begin{ex}
  गोलक चिन्हार्थमीमांसेत हे अनुमान
  अवैध ठरते; म्हणजे
  $p \cif r \Entails/
  (p \land q) \cif r$.
  $\mModel{M} = \tuple{W, O, V}$
  हे प्रतिमान पाहा: येथे
  $W = \{w, w_1, w_2\}$,
  $O_w = \{\{w\}, \{w,
  w_1\}, \{w, w_1, w_2\}\}$,
  $V(p) = \{w_1, w_2\}$,
  $V(q) = \{w_2\}$, आणि
  $V(r) = \{w_1\}$.
  इतर जगांसाठी $O_{w_1}=\{\{w_1\}\}$ आणि
  $O_{w_2}=\{\{w_2\}\}$ घ्या; त्यामुळे $O$ हे
  सर्व $W$ वर परिभाषित आहे.
  $p$ समाविष्ट करणारा
  $S = \{w, w_1\}$
  हा गोलक आहे आणि
  त्यातील सर्व जगांत
  $p \lif r$ सत्य आहे;
  म्हणून
  $\mSat{M}{p \cif r}[w]$.
  तसेच $(p \land q)$
  समाविष्ट करणारा
  $S' = \{w, w_1, w_2\}$
  हा गोलक आहे; पण
  $\mSat/{M}{(p \land q) \lif r}[w_2]$.
  $S'$ हा या $O_w$ मधील एकमेव $(p\land q)$-जग
  धारण करणारा गोलक आहे आणि त्यात अभिव्यंजन असत्य आहे.
  म्हणून कोणताही साक्षी-गोलक नाही; रिक्तपणे सत्य होण्याची
  अटही लागू होत नाही. त्यामुळे
  $\mSat/{M}{(p \land q) \cif r}[w]$
  (\ref{cnt:min:agg:fig:antecedent-strengthening} पाहा).

\begin{figure}
\begin{center}
\begin{tikzpicture}[layerwidth=1.5,scale=.6]\tiny
  \spheresystem{3}
  \propositionintersect{-10}{3}{30}{5}
  \begin{scope}
    \clip plot[smooth,tension=1] coordinates
          {(0,4.5) (30:.95) (4.5,-3.5)} -- (4.5,4.5) ;
    \spherefill{2}
  \end{scope}
  \draw[proposition] plot[dashed,smooth,tension=1] coordinates
          {(0,4.5) (30:.95) (4.5,-3.5)} ;
  \proposition{30}{2}{30}{5}
  \path (0,0) node[world] {$w$};
  \spherepos{30}{2}{node[world] {$w_1$}}
  \spherepos{-10}{3}{node[world] {$w_2$}}
  \spherepos{-10}{4}{node {$q$}}
  \spherepos{30}{4}{node {$r$}}
  \draw (1,4.5) node[below] {$p$};
\end{tikzpicture}
\caption{पूर्वांग बळकट करणे या नियमाचे प्रतिउदाहरण}
\label{cnt:min:agg:fig:antecedent-strengthening}
\end{center}
\end{figure}
\end{ex}












\section{संक्रमकता}\label{cnt:min:tra:sec}

वास्तविक अभिव्यंजनासाठी
साखळी-नियम लागू होतो:
$\varphi \lif \psi, \psi
\lif \chi \Entails \varphi \lif \chi$.
दुसऱ्या शब्दांत,
वास्तविक अभिव्यंजन
संक्रमक आहे.
प्रतिवास्तविक विधानांसाठीही
हे सत्य आहे का?
स्टालनाकर यांनी दिलेले
पुढील उदाहरण पाहा.
\begin{quote}
  जे.~एडगर हूव्हर रशियात
  जन्मले असते, तर ते
  कम्युनिस्ट झाले असते.

  जे.~एडगर हूव्हर कम्युनिस्ट
  असते, तर ते देशद्रोही
  झाले असते.

  म्हणून जे.~एडगर हूव्हर
  रशियात जन्मले असते,
  तर ते देशद्रोही
  झाले असते.
\end{quote}
हूव्हर यांचा जन्म
(प्रत्यक्षात ज्या वेळी झाला
त्याच वेळी) अमेरिकेत
न होता रशियात झाला
असता, तर ते
सोव्हिएत संघात वाढले
असते आणि कम्युनिस्ट
झाले असते, असे
समजू या. म्हणून
पहिले गृहीतक सत्य
आहे. त्याचप्रमाणे,
दुसरे गृहीतक
स्वतंत्रपणे पाहिल्यास
सत्य आहे. पण
निष्कर्ष असत्य आहे:
हूव्हर यांचा जन्म
सोव्हिएत संघात
झाला असता, तर
ते बहुधा निष्ठावान
कम्युनिस्ट झाले
असते; आपल्या
देशाशी विश्वासघात
केला नसता.
सत्य-मूल्यांची ही
सहज वाटणारी
नेमणूक स्टालनाकर--लुईस
यांच्या विश्लेषणातूनही
मिळते. हूव्हर यांच्या
जन्मस्थळातील बदल आणि त्याचे आवश्यक परिणाम
यांखेरीज अनावश्यक बदल नसलेले आपल्या जगाच्या
सर्वाधिक जवळचे
संभाव्य जग म्हणजे
हूव्हर सोव्हिएत
संघाचे चांगले नागरिक
म्हणून वाढतात
असे जग. पहिल्या
गृहीतकाचे आणि
निष्कर्षाचे पूर्वांग
सत्य असलेले हे
सर्वाधिक जवळचे
संभाव्य जग आहे.
तेथे हूव्हर
कम्युनिस्ट पक्षाचे
निष्ठावान सदस्य
असतात, म्हणून
देशद्रोही नसतात.
पण दुसरे गृहीतक
तपासताना वेगळे
जग पाहावे लागते:
हूव्हर कम्युनिस्ट
असलेले सर्वाधिक
जवळचे जग असे
की ते अमेरिकेत
जन्मले, नंतर
कम्युनिस्ट झाले,
आणि त्यामुळे
देशद्रोही ठरले.\footnote{
  अर्थात या उदाहरणाचा
  मुद्दा समजण्यासाठी
  काही सत्ताविषयक
  आणि राजकीय
  गृहीतके मानावी
  लागतात. उदाहरणार्थ,
  हूव्हर यांचा जन्म
  रशियन पालकांपासून
  होणे शक्य होते,
  किंवा १९५० च्या
  दशकातील अमेरिकेतील
  कम्युनिस्ट हे आपल्या
  देशाचे देशद्रोही होते,
  अशी गृहीतके.}

\begin{prob}
  प्रतिवास्तविक विधानांची
  संक्रमकता अपयशी
  ठरते याचे पटण्याजोगे,
  सहज उदाहरण शोधा.
\end{prob}

\begin{ex}\label{cnt:min:tra:ex:trans-counterex}
  गोलक चिन्हार्थमीमांसेत हे
  अनुमान अवैध ठरते;
  म्हणजे
  $p \cif q, q \cif r
  \Entails/ p \cif r$.
  $\mModel{M} = \tuple{W, O, V}$
  हे प्रतिमान पाहा: येथे
  $W = \{w, w_1, w_2\}$,
  $O_w = \{\{w\}, \{w,
  w_1\}, \{w, w_1, w_2\}\}$,
  $V(p) = \{w_2\}$,
  $V(q) = \{w_1, w_2\}$,
  आणि $V(r) = \{w_1\}$.
  तसेच $O_{w_1}=\{\{w_1\}\}$ आणि $O_{w_2}=\{\{w_2\}\}$
  घ्या, म्हणजे सर्व जगांचे गोलक-फलन पूर्ण होते.
  $p$ समाविष्ट करणारा
  $S = \{w, w_1, w_2\}$
  हा गोलक आहे आणि
  त्यातील सर्व जगांत
  $p \lif q$ सत्य आहे;
  म्हणून
  $\mSat{M}{p \cif q}[w]$.
  तसेच $q$ समाविष्ट
  करणारा
  $S' = \{w, w_1\}$
  हा गोलक आहे आणि
  त्यातील सर्व जगांत
  $\mSat{M}{q \lif r}$
  सत्य आहे; म्हणून
  $\mSat{M}{q \cif r}[w]$.
  पण $p$ समाविष्ट
  करणाऱ्या
  $\{w, w_1, w_2\}$
  गोलकात~$w_2$
  हे जग आहे,
  जिथे
  $\mSat/{M}{p \lif
    r}[w_2]$. हा $O_w$ मधील एकमेव $p$-जग धारण करणारा
  गोलक आहे; म्हणून $\mSat/{M}{p\cif r}[w]$.
  दोन्ही गृहीतके सत्य असून निष्कर्ष असत्य आहे.
\end{ex}

\begin{prob}
  \ref{cnt:min:tra:ex:trans-counterex}
  मधील प्रतिउदाहरणाला
  अनुरूप गोलकांची
  आकृती काढा.
\end{prob}

\begin{prob}
  \ref{cnt:min:tra:ex:trans-counterex}
  मध्ये जग~$w_2$
  येथे हूव्हर यांचा
  जन्म रशियात झाला
  आहे, ते कम्युनिस्ट
  आहेत आणि देशद्रोही
  नाहीत. जग~$w_1$
  येथे त्यांचा जन्म
  अमेरिकेत झाला
  आहे, ते कम्युनिस्ट
  आहेत आणि देशद्रोही
  आहेत. या प्रतिमानात
  $w_1$ हे~$w_2$
  पेक्षा~$w$ च्या
  अधिक जवळ आहे.
  हे आवश्यक आहे
  का? हूव्हर यांचा
  रशियात जन्म होणे
  हे त्यांच्या
  कम्युनिस्ट असण्यापेक्षा
  अधिक दूरची
  शक्यता आहे असे
  न मानणारे
  प्रतिउदाहरण
  देता येईल का?
\end{prob}












\section{प्रतिपरिवर्तन}\label{cnt:min:cpo:sec}

वास्तविक आणि कठोर
अभिव्यंजने त्यांच्या
प्रतिपरिवर्तनांशी सममूल्य
असतात. प्रतिवास्तविक
विधानांच्या बाबतीत
तसे नाही. क्रात्सर
यांनी दिलेले उदाहरण
पाहा:
\begin{quote}
  गटे यांचा मृत्यू १८३२ मध्ये
  झाला नसता, तरी आता
  ते मृतच असते.

  गटे आता मृत नसते,
  तर त्यांचा मृत्यू
  १८३२ मध्ये झाला असता.
\end{quote}
पहिले वाक्य सत्य आहे:
माणसे शेकडो वर्षे
जगत नाहीत. दुसरे
स्पष्टपणे असत्य आहे:
गटे आता मृत नसते,
तर ते अजून जिवंत
असते; त्यामुळे
त्यांचा मृत्यू
१८३२ मध्ये
झालेला नसता.

\begin{ex}\label{cnt:min:cpo:ex:contraposition-counterex}
  गोलक चिन्हार्थमीमांसेत
  प्रतिपरिवर्तन अवैध ठरते;
  म्हणजे
  $p \cif q \Entails/
  \lnot q \cif \lnot p$.
  $p$ चा अर्थ ‘‘गटे
  यांचा मृत्यू १८३२
  मध्ये झाला नाही’’
  आणि~$q$ चा अर्थ
  ‘‘गटे आता मृत
  आहेत’’ असा घ्या.
  हे पुढील प्रतिमानात
  दाखवता येते:
  $\mModel{M_1} = \tuple{W, O, V}$,
  येथे $W =
  \{w, w_1, w_2\}$,
  $O_w = \{\{w\}, \{w, w_1\},
  \{w, w_1, w_2\}\}$,
  $V(p) = \{w_1, w_2\}$
  आणि $V(q) = \{w, w_1\}$.
  $O_{w_1}=\{\{w_1\}\}$ आणि $O_{w_2}=\{\{w_2\}\}$ घ्या;
  $O$ आता संपूर्ण $W$ वर परिभाषित आहे.
  म्हणून~$w$ हे वास्तविक
  जग आहे, ज्यात गटे
  यांचा मृत्यू १८३२
  मध्ये झाला आणि
  ते अजूनही मृत आहेत;
  $w_1$ हे (जवळचे)
  जग आहे, ज्यात
  त्यांचा मृत्यू,
  समजा, १८३३ मध्ये
  झाला आणि ते
  अजूनही मृत आहेत;
  तर~$w_2$ हे
  (दूरचे) जग आहे,
  ज्यात गटे
  अजून जिवंत आहेत.
\begin{figure}
\begin{center}
\begin{tikzpicture}[scale=.6,layerwidth=1.5]\tiny
  \spheresystem{3}
  \begin{scope}
    \clip \propositionplot{0}{.8}{50}{4.5} -- (4.5,-4.5) -- (-4.5,-4.5) -- (-4.5,4.5) -- (4.5,4.5);
    \spherelayer{3}
  \end{scope}
  \propositionintersect{0}{2}{110}{5.5}
  \proposition{0}{.8}{50}{4.5}
  \path (0,0) node[world] {$w$};
  \spherepos{0}{2}{node[world] {$w_1$}}
  \spherepos{-50}{3}{node[world] {$w_2$}}
  \spherepos{28}{3}{node {$q$}}
  \spherepos{42}{3}{node {$\lnot q$}}
  \spherepos{-55}{4}{node {$p$}}
  \spherepos{-67}{4}{node {$\lnot p$}}
\end{tikzpicture}
\caption{प्रतिपरिवर्तनाचे प्रतिउदाहरण}
\label{cnt:min:cpo:fig:contraposition}
\end{center}
\end{figure}
$p$ समाविष्ट करणारा
$S = \{w, w_1\}$
हा गोलक आहे आणि
त्यातील सर्व जगांत
$p \lif q$ सत्य आहे;
म्हणून
$\mSat{M_1}{p \cif q}[w]$.
पण $\lnot q$
समाविष्ट करणाऱ्या
$\{w, w_1, w_2\}$
गोलकात~$w_2$
हे जग आहे; तेथे
$q$ असत्य आणि
$p$ सत्य आहे.
म्हणून
$\mSat/{M_1}{\lnot q
\lif \lnot p}[w_2]$. हा $O_w$ मधील एकमेव $\lnot q$-जग
धारण करणारा गोलक असल्याने
$\mSat/{M_1}{\lnot q\cif\lnot p}[w]$.
मूळ प्रतिवास्तविक विधान सत्य असून त्याचे प्रतिपरिवर्तन असत्य आहे.
\end{ex}



\part{संच उपपत्ती}\label{sth:::part}\label{sth:part}
